\chapter{Échantillonnage et fluctuation}

\section{Simulation d'échantillons}

\begin{defin}
	\nline
	\begin{itemize}
		\item Un \textbf{échantillon} de taille $n$ est constitué des résultats de $n$ répétitions indépendantes et identiquement distribuées d'une même expérience aléatoire.
		\item La \textbf{simulation} consiste à réaliser une expérience aléatoire à l'aide d'un algorithme ou d'un programme informatique générant des nombres aléatoires.
	\end{itemize}
\end{defin}

\begin{ex}
	On lance un dé à $6$ faces $n$ fois de suite et on observe le nombre de fois où le dé s'arrête sur la face \epsdice{1}.

	L'obtention d'un \epsdice{1} constitue le \og succès \fg{}, de probabilité théorique $p = \frac{1}{6} \approx 0{,}16666\ldots$.

	\begin{enumerate}
		\item \textbf{Génération d'un échantillon de lancers en Python :}

		      Le programme suivant renvoie la liste des résultats de $n$ lancers successifs :

		      \begin{lstlisting}
from random import randint

def echantillon(n):
    L = []
    for i in range(n):
        x = randint(1, 6)
        L.append(x)
    return L
\end{lstlisting}

		      Une exécution pour $n = 10$ donne par exemple :
		      \begin{lstlisting}
>>> echantillon(10)
[6, 4, 4, 1, 2, 6, 4, 5, 5, 6]
\end{lstlisting}

		\item \textbf{Calcul de la fréquence observée :}

		      On modifie la fonction afin qu'elle renvoie la fréquence d'apparition du \epsdice{1} sur $n$ lancers :

		      \begin{lstlisting}
from random import randint

def frequence_succes(n):
    c = 0
    for i in range(n):
        x = randint(1, 6)
        if x == 1:
            c = c + 1
    return c / n
\end{lstlisting}

		      \newpage

		      En augmentant la taille $n$ de l'échantillon, les fréquences observées se rapprochent de la valeur théorique $p = \frac{1}{6}$ :
		      \begin{lstlisting}
>>> frequence_succes(10)
0.1
>>> frequence_succes(100)
0.2
>>> frequence_succes(1000)
0.164
>>> frequence_succes(10000)
0.1613
>>> frequence_succes(10000000)
0.1665664
\end{lstlisting}
	\end{enumerate}
\end{ex}

\section{Fluctuation d'échantillonnage}

\begin{defin}
	La \textbf{fluctuation d'échantillonnage} est le phénomène par lequel la fréquence d'un événement varie d'un échantillon à un autre de même taille $n$, prélevés dans une même population.
\end{defin}

\begin{ex}
	Pour observer la fluctuation, on simule $N$ échantillons de même taille $n$ à l'aide de la fonction suivante :

	\begin{lstlisting}
from random import randint

def simulation(N, n):
    L = []
    for i in range(N):
        c = 0
        for j in range(n):
            if randint(1, 6) == 1:
                c = c + 1
        L.append(c / n)
    return L
\end{lstlisting}

	Pour $N = 10$ échantillons de taille $n = 50$, on obtient des fréquences fluctuantes :
	\begin{lstlisting}
>>> simulation(10, 50)
[0.3, 0.16, 0.12, 0.18, 0.22, 0.14, 0.14, 0.14, 0.24, 0.12]
\end{lstlisting}

	Pour $N = 10$ échantillons de taille $n = 10\,000$, la fluctuation est nettement plus faible et les fréquences sont très proches de $p \approx 0{,}166666\ldots$ :
	\begin{lstlisting}
>>> simulation(10, 10000)
[0.1706, 0.1632, 0.169, 0.1591, 0.1707, 0.1628, 0.169, 0.1684, 0.1627, 0.1656]
\end{lstlisting}
\end{ex}

\newpage

\begin{prop}
	Lorsque la taille $n$ de l'échantillon augmente, l'amplitude de la fluctuation d'échantillonnage diminue.

	Les fréquences observées se concentrent de plus en plus autour de la proportion théorique $p$.
\end{prop}

\begin{rem}
	Le nuage de points ci-dessous représente la simulation de $400$ échantillons de taille $n = 50$. Chaque point représente la fréquence observée $f$ du succès pour un échantillon donné.

	\begin{center}
		\begin{tikzpicture}[x=0.30mm, y=120mm]
			\draw[gray!30, ystep=0.05, xstep=50] (0,0) grid (400,0.4);
			\draw[thick, ->] (0,0) -- (420,0) node[right] {Échantillon};
			\draw[thick, ->] (0,0) -- (0,0.4) node[above] {Fréquence $f$};
			\foreach \x in {0,50,...,400} \draw (\x,0.005) -- (\x,-0.005) node[below] {\x};
			\foreach \y in {0.00,0.05,...,0.40} \draw (2,\y) -- (-2,\y) node[left] {\pgfmathprintnumber[fixed,precision=2]{\y}};
			\draw[rouge, dashed, very thick] (0,0.1666667) -- (400,0.1666667) node[right] {$p = \frac{1}{6}$};
			\draw[bleu] plot[only marks, mark=x, mark size=1.5pt] coordinates {
					(1,0.2) (2,0.18) (3,0.24) (4,0.16) (5,0.16) (6,0.18) (7,0.22) (8,0.16) (9,0.08) (10,0.2) (11,0.16) (12,0.14) (13,0.08) (14,0.28) (15,0.18) (16,0.2) (17,0.24) (18,0.22) (19,0.12) (20,0.1) (21,0.2) (22,0.26) (23,0.12) (24,0.1) (25,0.16) (26,0.2) (27,0.24) (28,0.18) (29,0.2) (30,0.2) (31,0.18) (32,0.22) (33,0.14) (34,0.26) (35,0.14) (36,0.1) (37,0.2) (38,0.18) (39,0.18) (40,0.16)
					(41,0.22) (42,0.18) (43,0.24) (44,0.26) (45,0.16) (46,0.18) (47,0.24) (48,0.16) (49,0.06) (50,0.16) (51,0.2) (52,0.18) (53,0.12) (54,0.18) (55,0.08) (56,0.12) (57,0.12) (58,0.1) (59,0.1) (60,0.2) (61,0.1) (62,0.2) (63,0.18) (64,0.18) (65,0.22) (66,0.18) (67,0.14) (68,0.24) (69,0.12) (70,0.2) (71,0.18) (72,0.26) (73,0.16) (74,0.22) (75,0.14) (76,0.22) (77,0.18) (78,0.12) (79,0.18) (80,0.1)
					(81,0.1) (82,0.16) (83,0.12) (84,0.2) (85,0.14) (86,0.2) (87,0.12) (88,0.14) (89,0.22) (90,0.12) (91,0.16) (92,0.2) (93,0.1) (94,0.1) (95,0.2) (96,0.22) (97,0.14) (98,0.18) (99,0.2) (100,0.22) (101,0.28) (102,0.1) (103,0.18) (104,0.24) (105,0.14) (106,0.16) (107,0.14) (108,0.2) (109,0.24) (110,0.22) (111,0.18) (112,0.14) (113,0.18) (114,0.14) (115,0.18) (116,0.14) (117,0.16) (118,0.16) (119,0.16) (120,0.28)
					(121,0.12) (122,0.2) (123,0.1) (124,0.2) (125,0.16) (126,0.18) (127,0.18) (128,0.16) (129,0.2) (130,0.14) (131,0.18) (132,0.14) (133,0.08) (134,0.18) (135,0.14) (136,0.08) (137,0.18) (138,0.24) (139,0.18) (140,0.18) (141,0.16) (142,0.22) (143,0.26) (144,0.16) (145,0.16) (146,0.2) (147,0.16) (148,0.16) (149,0.14) (150,0.2) (151,0.12) (152,0.16) (153,0.1) (154,0.08) (155,0.12) (156,0.12) (157,0.12) (158,0.12) (159,0.12) (160,0.14)
					(161,0.1) (162,0.14) (163,0.18) (164,0.1) (165,0.18) (166,0.16) (167,0.2) (168,0.18) (169,0.16) (170,0.24) (171,0.12) (172,0.16) (173,0.18) (174,0.18) (175,0.06) (176,0.16) (177,0.3) (178,0.2) (179,0.1) (180,0.22) (181,0.14) (182,0.16) (183,0.3) (184,0.18) (185,0.08) (186,0.1) (187,0.2) (188,0.18) (189,0.14) (190,0.08) (191,0.14) (192,0.2) (193,0.14) (194,0.2) (195,0.18) (196,0.14) (197,0.14) (198,0.24) (199,0.14) (200,0.18)
					(201,0.22) (202,0.12) (203,0.24) (204,0.1) (205,0.22) (206,0.14) (207,0.24) (208,0.28) (209,0.18) (210,0.1) (211,0.16) (212,0.16) (213,0.14) (214,0.24) (215,0.18) (216,0.3) (217,0.16) (218,0.2) (219,0.2) (220,0.1) (221,0.22) (222,0.28) (223,0.24) (224,0.18) (225,0.18) (226,0.1) (227,0.12) (228,0.22) (229,0.18) (230,0.18) (231,0.22) (232,0.16) (233,0.14) (234,0.18) (235,0.06) (236,0.14) (237,0.24) (238,0.2) (239,0.06) (240,0.16)
					(241,0.1) (242,0.06) (243,0.14) (244,0.24) (245,0.24) (246,0.16) (247,0.26) (248,0.18) (249,0.14) (250,0.22) (251,0.26) (252,0.2) (253,0.18) (254,0.22) (255,0.1) (256,0.08) (257,0.06) (258,0.08) (259,0.12) (260,0.1) (261,0.12) (262,0.16) (263,0.18) (264,0.14) (265,0.2) (266,0.28) (267,0.12) (268,0.16) (269,0.08) (270,0.22) (271,0.3) (272,0.16) (273,0.16) (274,0.16) (275,0.22) (276,0.28) (277,0.1) (278,0.16) (279,0.28) (280,0.14)
					(281,0.12) (282,0.12) (283,0.14) (284,0.36) (285,0.12) (286,0.08) (287,0.14) (288,0.18) (289,0.12) (290,0.08) (291,0.1) (292,0.12) (293,0.18) (294,0.14) (295,0.12) (296,0.16) (297,0.12) (298,0.14) (299,0.16) (300,0.1) (301,0.2) (302,0.12) (303,0.1) (304,0.1) (305,0.14) (306,0.16) (307,0.22) (308,0.12) (309,0.18) (310,0.28) (311,0.16) (312,0.1) (313,0.16) (314,0.16) (315,0.1) (316,0.14) (317,0.12) (318,0.22) (319,0.22) (320,0.12)
					(321,0.18) (322,0.16) (323,0.18) (324,0.22) (325,0.24) (326,0.16) (327,0.24) (328,0.22) (329,0.16) (330,0.22) (331,0.14) (332,0.18) (333,0.12) (334,0.06) (335,0.12) (336,0.18) (337,0.28) (338,0.18) (339,0.16) (340,0.12) (341,0.24) (342,0.1) (343,0.1) (344,0.2) (345,0.18) (346,0.18) (347,0.22) (348,0.2) (349,0.18) (350,0.18) (351,0.12) (352,0.32) (353,0.1) (354,0.2) (355,0.14) (356,0.1) (357,0.36) (358,0.14) (359,0.2) (360,0.08)
					(361,0.14) (362,0.18) (363,0.16) (364,0.16) (365,0.18) (366,0.18) (367,0.18) (368,0.14) (369,0.18) (370,0.14) (371,0.2) (372,0.26) (373,0.3) (374,0.18) (375,0.12) (376,0.28) (377,0.12) (378,0.1) (379,0.28) (380,0.2) (381,0.12) (382,0.2) (383,0.1) (384,0.12) (385,0.16) (386,0.04) (387,0.3) (388,0.18) (389,0.04) (390,0.12) (391,0.3) (392,0.12) (393,0.16) (394,0.16) (395,0.2) (396,0.12) (397,0.14) (398,0.18) (399,0.2) (400,0.14)
				};
		\end{tikzpicture}
	\end{center}

	Pour des échantillons de taille $n = 50$, les fréquences observées varient principalement entre $0{,}08$ et $0{,}30$.
\end{rem}

\section{Dispersion des résultats et intervalles}

\begin{prop}
	Soit un échantillon de taille $n$ issu d'une population où la proportion théorique d'un caractère est $p$.

	L'\textbf{écart-type} $s$ de la série des fréquences obtenues lors de répétitions d'échantillonnage peut être estimé par :
	\[
		s \approx \dfrac{1}{2\sqrt{n}}
	\]

	En moyenne, lors d'un grand nombre de simulations d'échantillons de même taille $n$ :
	\begin{itemize}
		\item environ \textbf{68\,\%} des fréquences observées appartiennent à l'intervalle $\intff{p - s}{p + s}$ ;
		\item environ \textbf{95\,\%} des fréquences observées appartiennent à l'intervalle $\intff{p - 2s}{p + 2s}$ ;
		\item environ \textbf{99\,\%} des fréquences observées appartiennent à l'intervalle $\intff{p - 3s}{p + 3s}$.
	\end{itemize}
\end{prop}

\begin{ex}
	Pour $n = 50$ et $p = \frac{1}{6} \approx 0{,}166666\ldots$, on calcule l'écart-type théorique :
	\[
		s \approx \dfrac{1}{2\sqrt{50}} \approx 0{,}071
	\]

	Les trois intervalles de dispersion associés sont :
	\begin{itemize}
		\item $\intff{p - s}{p + s} = \intff{0{,}096}{0{,}238}$ (contient environ $68\,\%$ des échantillons) ;
		\item $\intff{p - 2s}{p + 2s} = \intff{0{,}025}{0{,}309}$ (contient environ $95\,\%$ des échantillons) ;
		\item $\intff{p - 3s}{p + 3s} = \intff{0}{0{,}380}$ (contient environ $99\,\%$ des échantillons).
	\end{itemize}
\end{ex}

\newpage

\begin{rem}
	Représentation de la dispersion des $400$ échantillons dans les différents intervalles autour de $p$ :

	\begin{center}
		\begin{tikzpicture}[x=0.30mm, y=120mm]
			\draw[gray!30, ystep=0.05, xstep=50] (0,0) grid (400,0.4);
			\draw[thick, ->] (0,0) -- (420,0) node[right] {Échantillon};
			\draw[thick, ->] (0,0) -- (0,0.4) node[above] {Fréquence $f$};
			\foreach \x in {0,50,...,400} \draw (\x,0.005) -- (\x,-0.005) node[below] {\x};
			\foreach \y in {0.00,0.05,...,0.40} \draw (2,\y) -- (-2,\y) node[left] {\pgfmathprintnumber[fixed,precision=2]{\y}};
			\draw[rouge, dashed, very thick] (0,0.1666667) -- (400,0.1666667) node[right] {$p = \frac{1}{6}$};
			\draw[bleu] plot[only marks, mark=x, mark size=1.5pt] coordinates {
					(1,0.2) (2,0.18) (3,0.24) (4,0.16) (5,0.16) (6,0.18) (7,0.22) (8,0.16) (9,0.08) (10,0.2) (11,0.16) (12,0.14) (13,0.08) (14,0.28) (15,0.18) (16,0.2) (17,0.24) (18,0.22) (19,0.12) (20,0.1) (21,0.2) (22,0.26) (23,0.12) (24,0.1) (25,0.16) (26,0.2) (27,0.24) (28,0.18) (29,0.2) (30,0.2) (31,0.18) (32,0.22) (33,0.14) (34,0.26) (35,0.14) (36,0.1) (37,0.2) (38,0.18) (39,0.18) (40,0.16)
					(41,0.22) (42,0.18) (43,0.24) (44,0.26) (45,0.16) (46,0.18) (47,0.24) (48,0.16) (49,0.06) (50,0.16) (51,0.2) (52,0.18) (53,0.12) (54,0.18) (55,0.08) (56,0.12) (57,0.12) (58,0.1) (59,0.1) (60,0.2) (61,0.1) (62,0.2) (63,0.18) (64,0.18) (65,0.22) (66,0.18) (67,0.14) (68,0.24) (69,0.12) (70,0.2) (71,0.18) (72,0.26) (73,0.16) (74,0.22) (75,0.14) (76,0.22) (77,0.18) (78,0.12) (79,0.18) (80,0.1)
					(81,0.1) (82,0.16) (83,0.12) (84,0.2) (85,0.14) (86,0.2) (87,0.12) (88,0.14) (89,0.22) (90,0.12) (91,0.16) (92,0.2) (93,0.1) (94,0.1) (95,0.2) (96,0.22) (97,0.14) (98,0.18) (99,0.2) (100,0.22) (101,0.28) (102,0.1) (103,0.18) (104,0.24) (105,0.14) (106,0.16) (107,0.14) (108,0.2) (109,0.24) (110,0.22) (111,0.18) (112,0.14) (113,0.18) (114,0.14) (115,0.18) (116,0.14) (117,0.16) (118,0.16) (119,0.16) (120,0.28)
					(121,0.12) (122,0.2) (123,0.1) (124,0.2) (125,0.16) (126,0.18) (127,0.18) (128,0.16) (129,0.2) (130,0.14) (131,0.18) (132,0.14) (133,0.08) (134,0.18) (135,0.14) (136,0.08) (137,0.18) (138,0.24) (139,0.18) (140,0.18) (141,0.16) (142,0.22) (143,0.26) (144,0.16) (145,0.16) (146,0.2) (147,0.16) (148,0.16) (149,0.14) (150,0.2) (151,0.12) (152,0.16) (153,0.1) (154,0.08) (155,0.12) (156,0.12) (157,0.12) (158,0.12) (159,0.12) (160,0.14)
					(161,0.1) (162,0.14) (163,0.18) (164,0.1) (165,0.18) (166,0.16) (167,0.2) (168,0.18) (169,0.16) (170,0.24) (171,0.12) (172,0.16) (173,0.18) (174,0.18) (175,0.06) (176,0.16) (177,0.3) (178,0.2) (179,0.1) (180,0.22) (181,0.14) (182,0.16) (183,0.3) (184,0.18) (185,0.08) (186,0.1) (187,0.2) (188,0.18) (189,0.14) (190,0.08) (191,0.14) (192,0.2) (193,0.14) (194,0.2) (195,0.18) (196,0.14) (197,0.14) (198,0.24) (199,0.14) (200,0.18)
					(201,0.22) (202,0.12) (203,0.24) (204,0.1) (205,0.22) (206,0.14) (207,0.24) (208,0.28) (209,0.18) (210,0.1) (211,0.16) (212,0.16) (213,0.14) (214,0.24) (215,0.18) (216,0.3) (217,0.16) (218,0.2) (219,0.2) (220,0.1) (221,0.22) (222,0.28) (223,0.24) (224,0.18) (225,0.18) (226,0.1) (227,0.12) (228,0.22) (229,0.18) (230,0.18) (231,0.22) (232,0.16) (233,0.14) (234,0.18) (235,0.06) (236,0.14) (237,0.24) (238,0.2) (239,0.06) (240,0.16)
					(241,0.1) (242,0.06) (243,0.14) (244,0.24) (245,0.24) (246,0.16) (247,0.26) (248,0.18) (249,0.14) (250,0.22) (251,0.26) (252,0.2) (253,0.18) (254,0.22) (255,0.1) (256,0.08) (257,0.06) (258,0.08) (259,0.12) (260,0.1) (261,0.12) (262,0.16) (263,0.18) (264,0.14) (265,0.2) (266,0.28) (267,0.12) (268,0.16) (269,0.08) (270,0.22) (271,0.3) (272,0.16) (273,0.16) (274,0.16) (275,0.22) (276,0.28) (277,0.1) (278,0.16) (279,0.28) (280,0.14)
					(281,0.12) (282,0.12) (283,0.14) (284,0.36) (285,0.12) (286,0.08) (287,0.14) (288,0.18) (289,0.12) (290,0.08) (291,0.1) (292,0.12) (293,0.18) (294,0.14) (295,0.12) (296,0.16) (297,0.12) (298,0.14) (299,0.16) (300,0.1) (301,0.2) (302,0.12) (303,0.1) (304,0.1) (305,0.14) (306,0.16) (307,0.22) (308,0.12) (309,0.18) (310,0.28) (311,0.16) (312,0.1) (313,0.16) (314,0.16) (315,0.1) (316,0.14) (317,0.12) (318,0.22) (319,0.22) (320,0.12)
					(321,0.18) (322,0.16) (323,0.18) (324,0.22) (325,0.24) (326,0.16) (327,0.24) (328,0.22) (329,0.16) (330,0.22) (331,0.14) (332,0.18) (333,0.12) (334,0.06) (335,0.12) (336,0.18) (337,0.28) (338,0.18) (339,0.16) (340,0.12) (341,0.24) (342,0.1) (343,0.1) (344,0.2) (345,0.18) (346,0.18) (347,0.22) (348,0.2) (349,0.18) (350,0.18) (351,0.12) (352,0.32) (353,0.1) (354,0.2) (355,0.14) (356,0.1) (357,0.36) (358,0.14) (359,0.2) (360,0.08)
					(361,0.14) (362,0.18) (363,0.16) (364,0.16) (365,0.18) (366,0.18) (367,0.18) (368,0.14) (369,0.18) (370,0.14) (371,0.2) (372,0.26) (373,0.3) (374,0.18) (375,0.12) (376,0.28) (377,0.12) (378,0.1) (379,0.28) (380,0.2) (381,0.12) (382,0.2) (383,0.1) (384,0.12) (385,0.16) (386,0.04) (387,0.3) (388,0.18) (389,0.04) (390,0.12) (391,0.3) (392,0.12) (393,0.16) (394,0.16) (395,0.2) (396,0.12) (397,0.14) (398,0.18) (399,0.2) (400,0.14)
				};
			\draw[vert, line width=1.5pt] (0,0.096) rectangle (400,0.238);
			\node[vert, anchor=east, fill=white] at (395,0.24) {$68\,\%$};
			\draw[rouge, line width=1.5pt] (0,0.025) rectangle (400,0.309);
			\node[rouge, anchor=east, fill=white] at (395,0.31) {$95\,\%$};
		\end{tikzpicture}
	\end{center}
\end{rem}
